
Can neural networks learn permutation symmetries from data, or must such structure be built into the architecture? This work studies the question in weight-space learning, where models predict properties of neural networks from their weights. Using the Model Zoo dataset of approximately 42,000 CIFAR-10 CNNs, we compare permutation-equivariant Neural Functional Networks (NFN) against matched-capacity MLPs across data scales from 1,000 to 40,000 models.

At N=40,000, MLPs achieve high prediction accuracy ($R^{2}=0.986)$ but fail to develop permutation invariance (invariance score 0.63 versus NFN's 1.0). This demonstrates that task performance and representation quality are separable: MLPs learn dataset-specific position-accuracy correlations rather than semantic weight structure. At small scales, the gap is larger: NFN achieves $R^{2}=0.95$ at N=1,000 where MLP achieves $R^{2}=0.35$. NFN's predictions are invariant to weight permutation with deviations below 1.$19\times 10^{-7}$, confirming mathematical guarantees.

These findings falsify the hypothesis that sufficient data diversity teaches invariance. Permutation equivariance provides sample efficiency that data quantity cannot substitute.

